\chapter{Regularization}
\lecture{14}{28 Oct.}{}

\begin{prev}
    Overfitting happens with excessive power, stochastic and deterministic
    noise, and limited data. Of the remedies in the driving analogy, `start
    from a simple model' and the data-side ones were covered; `put the brakes',
    regularization, was left for later.
\end{prev}

This chapter builds the brakes. Stepping back from a large hypothesis set to a
smaller one is the same as adding a constraint on the weights. Softening that
constraint gives a family of sets \(\mathcal{H}(C)\), and solving the
constrained problem turns into minimizing an \emph{augmented error}, which for
linear regression has a closed form. VC theory explains why this works, and
the chapter ends with other choices of regularizer and the question of how
much to regularize.

\section{Regularized Hypothesis Set}\label{sec:reg-hyp-set}

\source{slides p.2--7}

\subsection{Regularization: The Magic}

\begin{center}
    \begin{tikzpicture}
        \begin{axis}[ovfplot, name=reg, title={`regularized fit'},
            title style={font=\small}]
            \ovfFiveData
            \ovfFiveTarget
            \ovfFiveGood
            \legend{Data, Target, Fit}
        \end{axis}
        \node at (6.15,1.5) {\(\Longleftarrow\)};
        \begin{axis}[ovfplot, at={(7.0cm,0)}, title={overfit},
            title style={font=\small}]
            \ovfFiveData
            \ovfFiveTarget
            \ovfFiveFit
            \legend{Data, Target, Fit}
        \end{axis}
    \end{tikzpicture}
\end{center}

The right plot is the overfit of \autoref{sec:what-overfit}: five examples
from a 2nd-order target, fitted exactly by a 4th-order polynomial. The left
plot is what regularization makes of the same data, a curve that stays close
to the target. The arrow points from right to left: regularization starts from
the overfit and pulls it back.

\begin{itemize}
    \item The idea is to `step back' from \(\mathcal{H}_{10}\) to
        \(\mathcal{H}_2\) in the nested chain of polynomial hypothesis sets,
        i.e.\ to move left in the picture below.
        \begin{center}
            \nestedH
        \end{center}
    \item The name has a history: it comes from function approximation for
        \textbf{ill-posed problems}, where many functions fit the given
        information equally well and an extra condition is needed to pick
        one.
\end{itemize}

\begin{question}
    How to step back?
\end{question}

\subsection{Stepping Back as Constraint}

\begin{center}
    \nestedH
\end{center}

Use the \(Q\)-th order polynomial transform for \(x \in \mathbb{R}\),
\[
    \bm{\Phi}_Q(x) = \left( 1, x, x^2, \ldots, x^Q \right),
\]
followed by linear regression, and from now on write the weights
\(\tilde{\bm{w}}\) of \(\mathcal{Z}\)-space simply as \(\bm{w}\). A hypothesis
in each set is then
\begin{align*}
    \text{hypothesis } \bm{w} \text{ in } \mathcal{H}_{10}&: \quad w_0 + w_1 x
    + w_2 x^2 + w_3 x^3 + \ldots + w_{10} x^{10} ,\\
    \text{hypothesis } \bm{w} \text{ in } \mathcal{H}_2&: \quad w_0 + w_1 x +
    w_2 x^2 .
\end{align*}
The second line is the first one with the last eight terms gone. That is,
\(\mathcal{H}_2 = \mathcal{H}_{10}\) AND `constraint that \(w_3 = w_4 = \ldots
= w_{10} = 0\)'.

\begin{moral}
    Step back \(=\) \textbf{constraint}.
\end{moral}

\subsection{Regression with Constraint}

Written as sets of weight vectors, both live in the same space
\(\mathbb{R}^{10+1}\):
\[
    \mathcal{H}_{10} \equiv \left\{ \bm{w} \in \mathbb{R}^{10+1} \right\} ,
    \qquad
    \mathcal{H}_2 \equiv \left\{ \bm{w} \in \mathbb{R}^{10+1} \text{ while }
    w_3 = w_4 = \ldots = w_{10} = 0 \right\} .
\]
Regression on each is the same minimization, with or without the constraint:
\begin{align*}
    \text{regression with } \mathcal{H}_{10}&: \quad \min_{\bm{w} \in
    \mathbb{R}^{10+1}} \Ein(\bm{w}) ,\\
    \text{regression with } \mathcal{H}_2&: \quad \min_{\bm{w} \in
    \mathbb{R}^{10+1}} \Ein(\bm{w}) \quad \text{s.t. } w_3 = w_4 = \ldots =
    w_{10} = 0 .
\end{align*}

\begin{moral}
    Step back \(=\) constrained optimization of \(\Ein\).
\end{moral}

\begin{question}
    Why don't you just use \(\bm{w} \in \mathbb{R}^{2+1}\)? :-)
\end{question}

For \(\mathcal{H}_2\) itself that would indeed be simpler. The point of
writing it as a constraint is that the constraint can now be changed, giving
sets between \(\mathcal{H}_2\) and \(\mathcal{H}_{10}\) that have no such
shortcut.

\subsection{Regression with Looser Constraint}

The first change: instead of insisting that the \emph{last} eight weights are
zero, ask only that \emph{some} eight are.
\[
    \mathcal{H}_2 \equiv \left\{ \bm{w} \in \mathbb{R}^{10+1} \text{ while }
    w_3 = \ldots = w_{10} = 0 \right\} ,
    \qquad
    \mathcal{H}_2' \equiv \left\{ \bm{w} \in \mathbb{R}^{10+1} \text{ while }
    \ge 8 \text{ of } w_q = 0 \right\} .
\]
Counting the nonzero weights with the Iverson bracket, regression becomes
\begin{align*}
    \text{regression with } \mathcal{H}_2&: \quad \min_{\bm{w} \in
    \mathbb{R}^{10+1}} \Ein(\bm{w}) \quad \text{s.t. } w_3 = \ldots = w_{10}
    = 0 ,\\
    \text{regression with } \mathcal{H}_2'&: \quad \min_{\bm{w} \in
    \mathbb{R}^{10+1}} \Ein(\bm{w}) \quad \text{s.t. } \sum_{q=0}^{10}
    \llbracket w_q \neq 0 \rrbracket \le 3 .
\end{align*}
\begin{itemize}
    \item \(\mathcal{H}_2'\) is more flexible than \(\mathcal{H}_2\):
        \(\mathcal{H}_2 \subset \mathcal{H}_2'\), since any three weights may
        be the free ones, not just \(w_0, w_1, w_2\).
    \item \(\mathcal{H}_2'\) is less risky than \(\mathcal{H}_{10}\):
        \(\mathcal{H}_2' \subset \mathcal{H}_{10}\), since it still uses at
        most three weights.
\end{itemize}

\begin{moral}
    Bad news for \vocab{sparse hypothesis set} \(\mathcal{H}_2'\):
    \textbf{NP-hard to solve} :-(
\end{moral}

The difficulty is the choice of which three weights to keep. There are
\(\binom{11}{3}\) choices here, and in general the count grows combinatorially,
just like the 0/1 error of the perceptron.

\subsection{Regression with Softer Constraint}

The second change replaces `how many weights are nonzero' by `how large the
weights are':
\[
    \mathcal{H}_2' \equiv \left\{ \bm{w} \in \mathbb{R}^{10+1} \text{ while }
    \ge 8 \text{ of } w_q = 0 \right\} ,
    \qquad
    \mathcal{H}(C) \equiv \left\{ \bm{w} \in \mathbb{R}^{10+1} \text{ while }
    \norm{\bm{w}}^2 \le C \right\} ,
\]
\begin{align*}
    \text{regression with } \mathcal{H}_2'&: \quad \min_{\bm{w} \in
    \mathbb{R}^{10+1}} \Ein(\bm{w}) \quad \text{s.t. } \sum_{q=0}^{10}
    \llbracket w_q \neq 0 \rrbracket \le 3 ,\\
    \text{regression with } \mathcal{H}(C)&: \quad \min_{\bm{w} \in
    \mathbb{R}^{10+1}} \Ein(\bm{w}) \quad \text{s.t. } \sum_{q=0}^{10} w_q^2
    \le C .
\end{align*}
\begin{itemize}
    \item \(\mathcal{H}(C)\) overlaps with \(\mathcal{H}_2'\) but is not exactly
        the same: both keep the weights `small', one by count and one by
        size.
    \item Over \(C \ge 0\), the sets form a soft and smooth structure:
        \[
            \mathcal{H}(0) \subset \mathcal{H}(1.126) \subset \ldots \subset
            \mathcal{H}(1126) \subset \ldots \subset \mathcal{H}(\infty) =
            \mathcal{H}_{10} .
        \]
        Unlike \(\mathcal{H}_0 \subset \mathcal{H}_1 \subset \cdots\), the
        size can now be tuned continuously by \(C\).
\end{itemize}

\begin{definition}[Regularized hypothesis]\label{def:wreg}
    The \vocab{regularized hypothesis} \(\wreg\) is the optimal solution from
    the \vocab{regularized hypothesis set} \(\mathcal{H}(C)\).
\end{definition}

\begin{exercise}[Fun Time]
    For \(Q \ge 1\), which of the following hypothesis (weight vector \(\bm{w}
    \in \mathbb{R}^{Q+1}\)) is not in the regularized hypothesis set
    \(\mathcal{H}(1)\)?
    \begin{enumerate}[(1)]
        \item \(\bm{w}^\top = [0, 0, \ldots, 0]\);
        \item \(\bm{w}^\top = [1, 0, \ldots, 0]\);
        \item \(\bm{w}^\top = [1, 1, \ldots, 1]\);
        \item \(\bm{w}^\top = \left[ \sqrt{\frac{1}{Q+1}}, \sqrt{\frac{1}{Q+1}},
            \ldots, \sqrt{\frac{1}{Q+1}} \right]\).
    \end{enumerate}
\end{exercise}

\begin{proof}[Reference Answer]
    (3). The squared length of \(\bm{w}\) in (3) is \(Q + 1\), which is not
    \(\le 1\).
\end{proof}

\section{Weight Decay Regularization}\label{sec:weight-decay}

\source{slides p.8--13}

\subsection{Matrix Form of Regularized Regression Problem}

The regularized regression problem, with \(\bm{z}_n = \bm{\Phi}_Q(x_n)\):
\[
    \min_{\bm{w} \in \mathbb{R}^{Q+1}} \Ein(\bm{w}) = \frac{1}{N}
    \underbrace{\sum_{n=1}^{N} \left( \bm{w}^\top \bm{z}_n - y_n
    \right)^2}_{(\mathrm{Z}\bm{w} - \bm{y})^\top (\mathrm{Z}\bm{w} - \bm{y})}
    \quad \text{s.t.} \quad \underbrace{\sum_{q=0}^{Q} w_q^2}_{\bm{w}^\top
    \bm{w}} \le C .
\]
\begin{itemize}
    \item The sum \(\sum_n \ldots\) equals \((\mathrm{Z}\bm{w} -
        \bm{y})^\top (\mathrm{Z}\bm{w} - \bm{y})\), the matrix form of linear
        regression in \autoref{sec:linreg-algo}, remember? :-)
    \item \(\bm{w}^\top \bm{w} \le C\) says the feasible \(\bm{w}\) lie within
        a radius-\(\sqrt{C}\) hypersphere centred at the origin.
\end{itemize}

\begin{question}
    How to solve the constrained optimization problem?
\end{question}

\subsection{The Lagrange Multiplier}

\[
    \min_{\bm{w} \in \mathbb{R}^{Q+1}} \Ein(\bm{w}) = \frac{1}{N} (\mathrm{Z}
    \bm{w} - \bm{y})^\top (\mathrm{Z}\bm{w} - \bm{y}) \quad \text{s.t.} \quad
    \bm{w}^\top \bm{w} \le C
\]

\begin{center}
    \regconstraint{l2}
\end{center}

The picture is the problem in weight space. \(\bm{w}_{\mathrm{lin}}\), the
unconstrained minimizer, lies outside the red sphere, so the best feasible
\(\bm{w}\) sits on the sphere's boundary. The blue ellipse is the contour of
\(\Ein\) through the current \(\bm{w}\).
\begin{itemize}
    \item The decreasing direction of \(\Ein\) is \(-\nabla \Ein(\bm{w})\), as
        in gradient descent (\autoref{sec:gd}), remember? :-)
    \item The normal vector of the boundary \(\bm{w}^\top \bm{w} = C\) at
        \(\bm{w}\) is \(\bm{w}\) itself (red arrow), pointing straight out of
        the sphere.
    \item If \(-\nabla \Ein(\bm{w})\) and \(\bm{w}\) are not parallel, then
        \(-\nabla \Ein(\bm{w})\) has a component along the boundary (green
        arrow). Moving that way decreases \(\Ein(\bm{w})\) without violating
        the constraint, so \(\bm{w}\) is not yet optimal.
    \item Hence at the optimal solution \(\wreg\) no such component is left:
        \(-\nabla \Ein(\wreg) \propto \wreg\).
\end{itemize}

\begin{moral}
    Want: find \vocab{Lagrange multiplier} \(\lambda > 0\) and \(\wreg\) such
    that \(\nabla \Ein(\wreg) + \frac{2\lambda}{N} \wreg = \bm{0}\).
\end{moral}

The proportionality constant is written as \(\frac{2\lambda}{N}\) only to make
the next formulas clean; any positive constant would do.

\begin{remark}
    Why must \(\lambda\) be positive? The picture assumes \(\bm{w}_{\mathrm{lin}}\)
    lies outside the sphere, so the constraint is active and \(\wreg\) is on
    the boundary. There, \(\Ein\) decreases \emph{outward}, towards
    \(\bm{w}_{\mathrm{lin}}\), which means \(-\nabla \Ein(\wreg)\) points the
    same way as \(\wreg\) and the constant is positive. If it pointed inward
    (\(\lambda < 0\)), moving inward would lower \(\Ein\) and stay feasible, so
    \(\wreg\) would not be optimal. The remaining case is \(\lambda = 0\): then
    \(\nabla \Ein(\wreg) = \bm{0}\), i.e.\ \(\wreg = \bm{w}_{\mathrm{lin}}\),
    which happens when \(\bm{w}_{\mathrm{lin}}\) is already inside the sphere
    and the constraint does nothing. That is why the next subsection allows
    `\(\lambda > 0\) or \(\lambda = 0\)'.
\end{remark}

\subsection{Augmented Error}

\begin{itemize}
    \item If an oracle tells you \(\lambda > 0\), then solving
        \[
            \nabla \Ein(\wreg) + \frac{2\lambda}{N} \wreg = \bm{0}
        \]
        means, with the gradient of linear regression from
        \autoref{sec:linreg-algo},
        \[
            \frac{2}{N} \left( \mathrm{Z}^\top \mathrm{Z} \wreg -
            \mathrm{Z}^\top \bm{y} \right) + \frac{2\lambda}{N} \wreg =
            \bm{0} .
        \]
    \item The optimal solution is
        \[
            \wreg \leftarrow (\mathrm{Z}^\top \mathrm{Z} + \lambda
            \mathrm{I})^{-1} \mathrm{Z}^\top \bm{y} ,
        \]
        called \vocab{ridge regression} in Statistics. For \(\lambda > 0\) the
        matrix \(\mathrm{Z}^\top \mathrm{Z} + \lambda \mathrm{I}\) is always
        invertible, unlike \(\mathrm{Z}^\top \mathrm{Z}\) alone.
\end{itemize}

The same equation can also be read the other way round: it is exactly the
condition `gradient \(= \bm{0}\)' of an unconstrained function.
\begin{itemize}
    \item If an oracle tells you \(\lambda > 0\), then solving
        \(\nabla \Ein(\wreg) + \frac{2\lambda}{N} \wreg = \bm{0}\) is
        equivalent to minimizing
        \[
            \underbrace{\Ein(\bm{w}) + \frac{\lambda}{N} \overbrace{\bm{w}^\top
            \bm{w}}^{\text{regularizer}}}_{\text{augmented error }
            \Eaug(\bm{w})} ,
        \]
        since the gradient of \(\frac{\lambda}{N} \bm{w}^\top \bm{w}\) is
        \(\frac{2\lambda}{N} \bm{w}\).
    \item So regularization can be done with the \vocab{augmented error}
        instead of the constrained \(\Ein\):
        \[
            \wreg \leftarrow \argmin_{\bm{w}} \Eaug(\bm{w}) \quad \text{for
            given } \lambda > 0 \text{ or } \lambda = 0 .
        \]
\end{itemize}

\begin{moral}
    Minimizing \textbf{unconstrained} \(\Eaug\) effectively minimizes
    \textbf{some \(C\)-constrained} \(\Ein\).
\end{moral}

The trade is that we no longer specify \(C\) but \(\lambda\).

\begin{remark}
    `\(\lambda\) corresponds to \(C\)' means that both problems have the same
    answer: the \(\bm{w}\) that minimizes \(\Eaug\) with this \(\lambda\) is also
    the \(\bm{w}\) that minimizes \(\Ein\) under \(\bm{w}^\top \bm{w} \le C\).
    Here is how to find that \(C\) for a given \(\lambda > 0\).
    \begin{enumerate}[(1)]
        \item \textbf{Solve the \(\lambda\) problem.} Let \(\wreg\) be the
            minimizer of \(\Eaug(\bm{w}) = \Ein(\bm{w}) + \frac{\lambda}{N}
            \bm{w}^\top \bm{w}\).
        \item \textbf{Choose \(C\) from its answer.} Set \(C = \wreg^\top
            \wreg\). Then \(\wreg\) itself satisfies the constraint \(\bm{w}^\top
            \bm{w} \le C\), with equality.
        \item \textbf{Compare with every other feasible \(\bm{w}\).} Take any
            \(\bm{w}\) with \(\bm{w}^\top \bm{w} \le C\). Since \(\wreg\)
            minimizes \(\Eaug\),
            \[
                \Ein(\bm{w}) + \frac{\lambda}{N} \bm{w}^\top \bm{w} \;\ge\;
                \Ein(\wreg) + \frac{\lambda}{N} \underbrace{\wreg^\top
                \wreg}_{= C} .
            \]
            Move the penalty term to the right:
            \[
                \Ein(\bm{w}) \;\ge\; \Ein(\wreg) + \frac{\lambda}{N}
                \underbrace{\left( C - \bm{w}^\top \bm{w} \right)}_{\ge 0
                \text{ since } \bm{w} \text{ is feasible}} \;\ge\;
                \Ein(\wreg) .
            \]
        \item \textbf{Conclude.} No feasible \(\bm{w}\) has a smaller \(\Ein\)
            than \(\wreg\), so \(\wreg\) is the solution of the
            \(C\)-constrained problem.
    \end{enumerate}
    So each \(\lambda\) comes with the \(C\) equal to the squared length of its
    own solution. A larger \(\lambda\) charges more for length, so the
    solution comes out shorter, and its \(C\) is smaller.
\end{remark}

\subsection{The Results}

\begin{center}
    \begin{adjustbox}{max width=\textwidth}
        \begin{tikzpicture}
            \begin{axis}[ovfplot, width=4.4cm, height=3.8cm,
                title={\(\lambda = 0\)}, title style={font=\small}]
                \ovfFiveData
                \ovfFiveTarget
                \ovfFiveFit
                \legend{Data, Target, Fit}
            \end{axis}
            \begin{axis}[ovfplot, width=4.4cm, height=3.8cm, at={(4.9cm,0)},
                title={\(\lambda = 0.0001\)}, title style={font=\small}]
                \ovfFiveData
                \ovfFiveTarget
                \regLambdaSmall
            \end{axis}
            \begin{axis}[ovfplot, width=4.4cm, height=3.8cm, at={(9.8cm,0)},
                title={\(\lambda = 0.01\)}, title style={font=\small}]
                \ovfFiveData
                \ovfFiveTarget
                \regLambdaMid
            \end{axis}
            \begin{axis}[ovfplot, width=4.4cm, height=3.8cm, at={(14.7cm,0)},
                title={\(\lambda = 1\)}, title style={font=\small}]
                \ovfFiveData
                \ovfFiveTarget
                \regLambdaOne
            \end{axis}
            \node[font=\small] at (1.75,-1.6) {overfitting};
            \foreach \x in {4.2,9.1,14.0}{\node at (\x,-1.6) {\(\Longrightarrow\)};}
            \node[font=\small] at (16.45,-1.6) {underfitting};
        \end{tikzpicture}
    \end{adjustbox}
\end{center}

The same five examples, fitted with a 4th-order transform and increasing
\(\lambda\). At \(\lambda = 0\) there is no regularization and the fit is the
overfit from before. Already \(\lambda = 0.0001\) removes the large swing and
gives the `regularized fit' of the first section. At \(\lambda = 1\) the
weights are pressed so short that the fit can no longer follow the data, which
is underfitting.

\begin{moral}
    Philosophy: \textbf{a little regularization} goes a long way!
\end{moral}

\begin{definition}[Weight-decay regularization]\label{def:weight-decay}
    Adding \(+\frac{\lambda}{N} \bm{w}^\top \bm{w}\) to \(\Ein\) is called
    \vocab{weight-decay regularization}:
    \[
        \text{larger } \lambda \iff \text{prefer shorter } \bm{w} \iff
        \text{effectively smaller } C .
    \]
    It goes with `any' transform \(+\) linear model.
\end{definition}

The name comes from gradient descent: each step on \(\Eaug\) first shrinks the
current \(\bm{w}\) by a constant factor and then moves along \(-\nabla \Ein\),
so without support from the data the weights decay towards zero.

\subsection{Some Detail: Legendre Polynomials}

\[
    \min_{\bm{w} \in \mathbb{R}^{Q+1}} \frac{1}{N} \sum_{n=1}^{N} \left(
    \bm{w}^\top \bm{\Phi}(x_n) - y_n \right)^2 + \frac{\lambda}{N}
    \sum_{q=0}^{Q} w_q^2
\]
Weight decay penalizes every \(w_q\) by the same amount, so the transform
matters.
\begin{itemize}
    \item The naive polynomial transform is \(\bm{\Phi}(x) = \left( 1, x, x^2,
        \ldots, x^Q \right)\). When \(x_n \in [-1, +1]\), \(x_n^q\) is really
        small for large \(q\), so a high-order term needs a large \(w_q\) to
        have any effect, and weight decay punishes it unfairly.
    \item The normalized polynomial transform is \(\left( 1, L_1(x), L_2(x),
        \ldots, L_Q(x) \right)\), using the `orthonormal basis functions' called
        \vocab{Legendre polynomials}. Each \(L_q\) spans the whole range
        \([-1, 1]\), so equal weights mean comparable effects.
\end{itemize}

\begin{center}
    \begin{adjustbox}{max width=\textwidth}
        \begin{tikzpicture}
            \begin{groupplot}[group style={group size=5 by 1,
                    horizontal sep=0.25cm},
                width=3.6cm, height=3.0cm, xmin=-1, xmax=1, ymin=-1.1,
                ymax=1.1, xtick=\empty, ytick=\empty,
                title style={font=\small}, xlabel style={font=\scriptsize},
                domain=-1:1, samples=80]
                \nextgroupplot[title={\(L_1\)}, xlabel={\(x\)}]
                \addplot[blue!85!black, thick] {x};
                \nextgroupplot[title={\(L_2\)},
                    xlabel={\(\frac{1}{2}(3x^2 - 1)\)}]
                \addplot[blue!85!black, thick] {(3*x^2 - 1)/2};
                \nextgroupplot[title={\(L_3\)},
                    xlabel={\(\frac{1}{2}(5x^3 - 3x)\)}]
                \addplot[blue!85!black, thick] {(5*x^3 - 3*x)/2};
                \nextgroupplot[title={\(L_4\)},
                    xlabel={\(\frac{1}{8}(35x^4 - 30x^2 + 3)\)}]
                \addplot[blue!85!black, thick] {(35*x^4 - 30*x^2 + 3)/8};
                \nextgroupplot[title={\(L_5\)},
                    xlabel={\(\frac{1}{8}(63x^5 \cdots)\)}]
                \addplot[blue!85!black, thick]
                    {(63*x^5 - 70*x^3 + 15*x)/8};
            \end{groupplot}
        \end{tikzpicture}
    \end{adjustbox}
\end{center}

The plots show why: every \(L_q\) swings between \(-1\) and \(1\) on the whole
interval, with \(q\) zero crossings, whereas \(x^q\) stays near \(0\) except
close to \(\pm 1\).

\begin{exercise}[Fun Time]
    When would \(\wreg\) equal \(\bm{w}_{\mathrm{LIN}}\)?
    \begin{enumerate}[(1)]
        \item \(\lambda = 0\);
        \item \(C = \infty\);
        \item \(C \ge \norm{\bm{w}_{\mathrm{LIN}}}^2\);
        \item all of the above.
    \end{enumerate}
\end{exercise}

\begin{proof}[Reference Answer]
    (4). (1) and (2) shall be easy; (3) means that there are effectively no
    constraint on \(\bm{w}\), hence the equivalence: \(\bm{w}_{\mathrm{LIN}}\)
    already lies inside the sphere, so the constraint never binds.
\end{proof}

\section{Regularization and VC Theory}\label{sec:reg-vc}

\source{slides p.14--17}

\subsection{Regularization and VC Theory}

\begin{center}
    \begin{tikzpicture}[font=\small]
        \node[mlbox, text width=5.4cm] (con) at (0,0)
            {Regularization by Constrained-Minimizing \(\Ein\)\\[2pt]
            \(\min_{\bm{w}} \Ein(\bm{w})\) s.t.\ \(\bm{w}^\top \bm{w} \le C\)};
        \node[mlbox, text width=5.4cm] (vc) at (7.6,0)
            {VC Guarantee of Constrained-Minimizing \(\Ein\)\\[2pt]
            \(\Eout(\bm{w}) \le \Ein(\bm{w}) + \Omega(\mathcal{H}(C))\)};
        \node[mlbox, text width=5.4cm] (aug) at (0,-2.8)
            {Regularization by Minimizing \(\Eaug\)\\[2pt]
            \(\min_{\bm{w}} \Eaug(\bm{w}) = \Ein(\bm{w}) + \frac{\lambda}{N}
            \bm{w}^\top \bm{w}\)};
        \draw[mlarrow] (con) -- (vc);
        \draw[mlarrow, {Stealth[length=2.2mm]}-{Stealth[length=2.2mm]}]
            (con) -- (aug) node[midway, right] {\(C\) equivalent to some
            \(\lambda\)};
    \end{tikzpicture}
\end{center}

Read the diagram as a chain. Minimizing \(\Ein\) under \(\bm{w}^\top \bm{w} \le
C\) is learning with the hypothesis set \(\mathcal{H}(C)\), so the VC bound
applies with the complexity of \(\mathcal{H}(C)\) (top arrow). Minimizing
\(\Eaug\) gives the same solution as the constrained problem for some \(C\)
(vertical arrow). So minimizing \(\Eaug\) inherits that guarantee, even though
it searches all of \(\mathbb{R}^{Q+1}\).

\begin{moral}
    Minimizing \(\Eaug\): indirectly getting VC guarantee \textbf{without
    confining to \(\mathcal{H}(C)\)}.
\end{moral}

\subsection{Another View of Augmented Error}

Put the two formulas side by side:
\[
    \text{augmented error: } \Eaug(\bm{w}) = \Ein(\bm{w}) + \frac{\lambda}{N}
    \bm{w}^\top \bm{w} ,
    \qquad
    \text{VC bound: } \Eout(\bm{w}) \le \Ein(\bm{w}) + \Omega(\mathcal{H}) .
\]
Both add a complexity term to \(\Ein\), but of different things.
\begin{itemize}
    \item The regularizer \(\bm{w}^\top \bm{w}\), written \(\Omega(\bm{w})\),
        measures the complexity of a \emph{single hypothesis}.
    \item The generalization price \(\Omega(\mathcal{H})\) measures the
        complexity of a \emph{hypothesis set}.
    \item If \(\frac{\lambda}{N} \Omega(\bm{w})\) `represents'
        \(\Omega(\mathcal{H})\) well, then \(\Eaug\) is a better proxy of
        \(\Eout\) than \(\Ein\) is.
\end{itemize}

\begin{moral}
    Minimizing \(\Eaug\):\\
    (heuristically) operating with the better proxy;\\
    (technically) enjoying flexibility of whole \(\mathcal{H}\).
\end{moral}

\subsection{Effective VC Dimension}

\[
    \min_{\bm{w} \in \mathbb{R}^{\tilde{d}+1}} \Eaug(\bm{w}) = \Ein(\bm{w}) +
    \frac{\lambda}{N} \Omega(\bm{w})
\]
What is the model complexity of this learning?
\begin{itemize}
    \item Counted by the hypothesis set, \(\dvc(\mathcal{H}) = \tilde{d} + 1\),
        because all \(\bm{w}\) are `considered' during the minimization.
    \item The \(\bm{w}\) that are `actually needed', however, form only
        \(\mathcal{H}(C)\), with some \(C\) equivalent to \(\lambda\).
    \item \(\dvc(\mathcal{H}(C))\) is therefore the more honest count. It
        depends on the algorithm as well as on \(\mathcal{H}\), and is called
        the \vocab{effective VC dimension} \(\deff(\mathcal{H},
        \mathcal{A})\), where \(\mathcal{A}\) is `\(\min \Eaug\)'.
\end{itemize}

\begin{moral}
    Explanation of regularization: \(\dvc(\mathcal{H})\) large, while
    \(\deff(\mathcal{H}, \mathcal{A})\) small if \(\mathcal{A}\) regularized.
\end{moral}

\begin{exercise}[Fun Time]
    Consider the weight-decay regularization with regression. When increasing
    \(\lambda\) in \(\mathcal{A}\), what would happen with \(\deff(\mathcal{H},
    \mathcal{A})\)?
    \begin{enumerate}[(1)]
        \item \(\deff \uparrow\);
        \item \(\deff \downarrow\);
        \item \(\deff = \dvc(\mathcal{H})\) and does not depend on \(\lambda\);
        \item \(\deff = 1126\) and does not depend on \(\lambda\).
    \end{enumerate}
\end{exercise}

\begin{proof}[Reference Answer]
    (2). Larger \(\lambda \iff\) smaller \(C \iff\) smaller \(\mathcal{H}(C)
    \iff\) smaller \(\deff\).
\end{proof}

\section{General Regularizers}\label{sec:general-reg}

\source{slides p.18--21}

\subsection{General Regularizers \texorpdfstring{\(\Omega(\bm{w})\)}{Omega(w)}}

Weight decay is one choice of \(\Omega(\bm{w})\). In general we want the
constraint to point in the `direction' of the target function, so that what
it removes is mostly what we would have got wrong anyway.
\begin{itemize}
    \item \emph{Target-dependent}: use some properties of the target, if known.
        \begin{itemize}
            \item For a symmetric target, the \textbf{symmetry} regularizer
                \(\sum \llbracket q \text{ is odd} \rrbracket w_q^2\) penalizes
                only the odd terms.
        \end{itemize}
    \item \emph{Plausible}: push in a direction towards \textbf{smoother} or
        \textbf{simpler} hypotheses, because stochastic and deterministic
        noise are both \textbf{non-smooth}, and fitting them is what
        overfitting does.
        \begin{itemize}
            \item The \textbf{sparsity} (L1) regularizer \(\sum \abs{w_q}\)
                is one such choice (next subsection).
        \end{itemize}
    \item \emph{Friendly}: make the problem easy to optimize.
        \begin{itemize}
            \item The \textbf{weight-decay} (L2) regularizer \(\sum w_q^2\)
                keeps a closed-form solution.
        \end{itemize}
    \item \emph{Bad?} :-) No worries: a poor regularizer is guarded by
        \(\lambda\), since \(\lambda = 0\) switches it off.
\end{itemize}

\begin{moral}
    Augmented error \(=\) error \(\errhat\) \(+\) regularizer \(\Omega\).\\
    Regularizer: target-dependent, plausible, or friendly.\\
    Ringing a bell? :-)\\
    Error measure: user-dependent, plausible, or friendly.
\end{moral}

The bell is \autoref{def:errhat}: the algorithmic error measure was chosen by
the same three criteria. Both parts of \(\Eaug\) are design choices of the
algorithm, made for the same kinds of reasons.

\subsection{L2 and L1 Regularizer}

\begin{center}
    \regconstraint{l2}
    \hspace{1cm}
    \regconstraint{l1}
\end{center}

\begin{center}
    \begin{tabular}{p{0.42\textwidth}p{0.42\textwidth}}
        \textbf{L2 Regularizer} & \textbf{L1 Regularizer}\\[2pt]
        \(\displaystyle \Omega(\bm{w}) = \sum_{q=0}^{Q} w_q^2 =
        \norm{\bm{w}}_2^2\)
        &
        \(\displaystyle \Omega(\bm{w}) = \sum_{q=0}^{Q} \abs{w_q} =
        \norm{\bm{w}}_1\)\\[10pt]
        \begin{itemize}
            \item convex, differentiable everywhere;
            \item easy to optimize.
        \end{itemize}
        &
        \begin{itemize}
            \item convex, \textbf{not} differentiable everywhere;
            \item \textbf{sparsity} in solution.
        \end{itemize}
    \end{tabular}
\end{center}

The left picture is the one from the Lagrange multiplier. On the right the
sphere becomes the diamond \(\norm{\bm{w}}_1 = C\). Its normal direction on
each flat face is the \textbf{sign} vector of \(\bm{w}\), the same at every
point of that face. So along a face, \(-\nabla \Ein\) keeps a component that
slides \(\bm{w}\) further, until it reaches a corner. The corners lie on the
axes, where some \(w_q = 0\): the solution is often sparse.

\begin{moral}
    L1 useful if needing \textbf{sparse solution}.
\end{moral}

\subsection{The Optimal \texorpdfstring{\(\lambda\)}{lambda}}

\begin{center}
    \begin{tikzpicture}
        \begin{axis}[width=6.4cm, height=4.6cm, xmin=0, xmax=2, ymin=0,
            ymax=1.01, axis lines=left, axis line style={-},
            xtick={0.5,1,1.5,2}, ytick={0.25,0.5,0.75,1},
            tick label style={font=\scriptsize}, label style={font=\small},
            xlabel={Regularization Parameter, \(\lambda\)},
            ylabel={Expected \(\Eout\)}, title={stochastic noise},
            title style={font=\small}, clip=false]
            \regStochHigh \regStochMid \regStochZero
            \regStochHighMin \regStochMidMin \regStochZeroMin
            \node[font=\scriptsize, text=red!85!black, anchor=south] at
                (axis cs:1.65,0.79) {\(\sigma^2 = 0.5\)};
            \node[font=\scriptsize, text=green!60!black, anchor=south] at
                (axis cs:1.65,0.52) {\(\sigma^2 = 0.25\)};
            \node[font=\scriptsize, text=blue!85!black, anchor=south] at
                (axis cs:1.65,0.18) {\(\sigma^2 = 0\)};
        \end{axis}
        \begin{axis}[at={(7.2cm,0)}, width=6.4cm, height=4.6cm, xmin=0,
            xmax=2, ymin=0, ymax=0.6, axis lines=left, axis line style={-},
            xtick={0.5,1,1.5,2}, ytick={0.2,0.4,0.6},
            tick label style={font=\scriptsize}, label style={font=\small},
            xlabel={Regularization Parameter, \(\lambda\)},
            ylabel={Expected \(\Eout\)}, title={deterministic noise},
            title style={font=\small}, clip=false]
            \regDetHigh \regDetMid \regDetLow
            \regDetHighMin \regDetMidMin \regDetLowMin
            \node[font=\scriptsize, text=red!85!black, anchor=south] at
                (axis cs:1.65,0.49) {\(Q_f = 100\)};
            \node[font=\scriptsize, text=green!60!black, anchor=south] at
                (axis cs:1.65,0.36) {\(Q_f = 30\)};
            \node[font=\scriptsize, text=blue!85!black, anchor=south] at
                (axis cs:1.65,0.16) {\(Q_f = 15\)};
        \end{axis}
    \end{tikzpicture}
\end{center}

Each curve is the expected \(\Eout\) of weight decay as \(\lambda\) grows, and
the dot marks its minimum, the best \(\lambda\) for that setting. On the left
the stochastic noise level \(\sigma^2\) varies; on the right the target
complexity \(Q_f\), i.e.\ the deterministic noise.
\begin{itemize}
    \item More noise \(\iff\) more regularization needed: the dot moves right
        as \(\sigma^2\) or \(Q_f\) grows, and without noise (\(\sigma^2 = 0\))
        the best \(\lambda\) is \(0\). In driving terms, a more bumpy road
        \(\iff\) putting on the brakes more.
    \item The noise is unknown in practice, so it is important to make proper
        choices of \(\lambda\). Too little leaves the overfitting; too much
        climbs the right side of each curve.
\end{itemize}

\begin{question}
    How to choose? Stay tuned for the next lecture! :-)
\end{question}

\begin{exercise}[Fun Time]
    Consider using a regularizer \(\Omega(\bm{w}) = \sum_{q=0}^{Q} 2^q w_q^2\)
    to work with Legendre polynomial regression. Which kind of hypothesis
    does the regularizer prefer?
    \begin{enumerate}[(1)]
        \item symmetric polynomials satisfying \(h(x) = h(-x)\);
        \item low-dimensional polynomials;
        \item high-dimensional polynomials;
        \item no specific preference.
    \end{enumerate}
\end{exercise}

\begin{proof}[Reference Answer]
    (2). There is a higher `penalty' for higher-order terms, and hence the
    regularizer prefers low-dimensional polynomials.
\end{proof}

\hrulebar

\subsection*{Summary}

\begin{description}
    \item[Regularized Hypothesis Set] original \(\mathcal{H}\) \(+\)
        constraint.
    \item[Weight Decay Regularization] add \(\frac{\lambda}{N} \bm{w}^\top
        \bm{w}\) in \(\Eaug\).
    \item[Regularization and VC Theory] regularization decreases \(\deff\).
    \item[General Regularizers] target-dependent, [plausible], or [friendly].
\end{description}

Regularization trades a little \(\Ein\) for a large drop in effective
complexity, and its strength is set by a single number \(\lambda\). The
optimal-\(\lambda\) curves show that the right value depends on noise we
cannot see. Next: choosing from the so-many models and parameters.
